% 第 3 章 §3.3.6--§3.3.9（印刷页 133--137，PDF 143--147）
% 承接 ch03c：原文 3.3.1--3.3.7 已由该文件排出，此处从注 3.3.8 起，故先把共享计数器置为 7。
\setcounter{theorem}{7}

\subsection{二次映射 \texorpdfstring{$\pi_2\to\pi_3$}{π2 到 π3}}

对只有 $\pi_2$、$\pi_3$ 两个同伦群的空间，其 Postnikov $k$-不变量 $k_3$ 落在 $H^4(K(\pi_2(X),\pi_3(X)))$ 中；下面给出它的一种解释。对任意空间 $X$，关于 $\pi_2(X)$ 及更高同伦群交换性的标准论证，掩盖了同伦群上的高阶运算；其中第一个是 \emph{Whitehead 括号}（Whitehead bracket），即一族双线性映射 $\pi_m(X)\times\pi_n(X)\to\pi_{m+n-1}(X)$。要看清这些映射，最好利用 $X$ 的带基点环路空间 $\Omega X$ 以及同构 $\pi_m(X)\cong\pi_{m-1}(\Omega X)$。具体地说，换位子映射 $\Omega X\times\Omega X\to\Omega X$ 可以形变，从而把子空间 $\Omega X\times\{1\}\cup\{1\}\times\Omega X$ 压到恒等元上。给定表示 $\pi_{m,n}(X)$ 中类的 $\alpha\colon S^{m-1}\to\Omega X$、$\beta\colon S^{n-1}\to\Omega X$，复合
\[
\alpha\times\beta\colon S^{m-1}\times S^{n-1}\to\Omega X\times\Omega X\xrightarrow{\;[\,,\,]\;}\Omega X
\]
被压扁，从而通过
\[
S^{m-1}\times S^{n-1}\big/\bigl(S^{m-1}\times\{*\}\cup\{*\}\times S^{n-1}\bigr)\cong S^{m+n-2}
\]
分解；这就给出 $\pi_{m+n-2}(\Omega X)\cong\pi_{m+n-1}(X)$ 中的一个类。

\begin{annote}底本此段有两处记号需说明，译文均照排。其一，$H^4(K(\pi_2(X),\pi_3(X)))$ 中 Eilenberg--MacLane 空间的第二个参数本应是阶数；据同页末段 $k_3\in H^4(K(\pi_2,2);\pi_3)$，此处应为 $K(\pi_2(X),2)$。其二，$\pi_{m,n}(X)$ 在底本中未加定义；按上下文，$\alpha$、$\beta$ 分别表示 $\pi_m(X)$、$\pi_n(X)$ 中的类。[来源：上下文推断]\end{annote}

\begin{remark}当 $m=n=1$ 时，它化为 $\pi_1(X)$ 上的换位子括号。\end{remark}

双线性的 Whitehead 括号 $\pi_2\times\pi_2\to\pi_3$ 有一个二次加细（quadratic refinement）。这是因为可以算出：$\pi_3(S^2)$ 的生成元——由著名的 Hopf 纤维化表示——是 $1\in\pi_2(S^2)$ 的 Whitehead 平方的\emph{一半}。对任意空间 $X$，把任意 $\alpha\colon S^2\to X$ 与 Hopf 映射 $S^3\to S^2$ 预先复合，便得到二次映射 $\pi_2(X)\to\pi_3(X)$。

\begin{theorem}只有两个非平凡同伦群 $\pi_2(X)$ 与 $\pi_3(X)$ 的空间 $X$，在同伦意义下完全由这两个群以及二次映射 $\pi_2\to\pi_3$ 决定。\end{theorem}

同伦论告诉我们，该空间在同伦意义下完全由 Postnikov 不变量 $k_3\in H^4(K(\pi_2,2);\pi_3)$ 决定。Eilenberg 与 MacLane 的一个结果断言：这样的上同调类恰好对应于二次映射 $\pi_2\to\pi_3$；我们刚刚在一个方向上构造了这一对应。

\subsection{由 \texorpdfstring{$X$}{X} 得到的一个编织张量范畴}

现在让我们把 $k_3$ 更严格地融入一个编织张量范畴，事实上融入一个带状范畴。

\begin{definition}一个\emph{张量范畴}（tensor category）$\mathcal{T}$ 是一个（$\C$-）线性范畴，带有双线性的乘法函子 $m\colon\mathcal{T}\times\mathcal{T}\to\mathcal{T}$，含单位对象 $1$ 使 $m(1,x)=x=m(x,1)$，并带结合子 $\alpha_{x,y,z}\colon m(m(x,y),z)\to m(x,m(y,z))$，后者让我们在乘法中「移动括号」。结合子对四个对象满足一条融贯性恒等式（Stasheff 的五边形恒等式）。\footnote{MacLane 的融贯性定理保证，我们可以把 $\mathcal{T}$ 换成一个等价的范畴，在其中 $\alpha\equiv1$；但这可能需要把范畴扩大为一个对象多得多的等价范畴，或者可能破坏某些附加结构，例如连续性。}\end{definition}

为简单起见，我们已把单位同构取为严格的。五边形恒等式是陈述「括号移动次序的选择无关紧要」的一种精确方式。乘积 $m(x,y)$ 常记作 $x\otimes y$。

粗略地说，一个张量范畴若交换到一阶，就称为\emph{编织的}（braided）。记 $\tau$ 为范畴的平方 $\mathcal{B}\times\mathcal{B}$ 上的转置自同构。

\begin{definition}一个\emph{编织张量范畴}（braided tensor category）$\mathcal{B}$ 是这样一个张量范畴：它配备函子的一个双乘性编织同构 $\beta\colon m\xrightarrow{\sim}m\circ\tau$。也就是说，对每一对 $x,y$，给定同构 $\beta(x,y)\colon x\otimes y\to y\otimes x$，它关于 $x,y$ 是函子性的。若 $\beta^2=\Id$，则称该范畴为\emph{对称的}。\end{definition}

双乘性是指：对任意三个对象 $x,y,z$，下列图表交换：
\begin{center}
\begin{tikzcd}[row sep=large, column sep=large]
x\otimes y\otimes z \arrow[rr, "{\beta(x,y\otimes z)}"] \arrow[dr, "{\beta(x,y)\otimes\Id_z}"'] & & y\otimes z\otimes x \\
& y\otimes x\otimes z \arrow[ur, "{\Id_y\otimes\beta(x,z)}"'] &
\end{tikzcd}
\end{center}
\begin{center}
\begin{tikzcd}[row sep=large, column sep=large]
x\otimes y\otimes z \arrow[rr, "{\beta(x\otimes y,z)}"] \arrow[dr, "{\Id_x\otimes\beta(y,z)}"'] & & z\otimes x\otimes y \\
& x\otimes z\otimes y \arrow[ur, "{\beta(x,z)\otimes\Id_y}"'] &
\end{tikzcd}
\end{center}

我们省略了结合子 $\alpha$（把它们都取为 $\Id$）。

\begin{remark}若把结合子搁在一边，那么编码编织结构的一种方式是：指定每个 $n$ 重张量积函子 $\mathcal{T}^{\otimes n}\to\mathcal{T}$（由逐次乘法定义）都带有 $n$ 股辫群（braid group）$B_n$ 的一个自然交织作用，并与包含 $B_m\times B_n\to B_{m+n}$ 相容。在对称情形，这些作用通过对称群分解。\end{remark}

\begin{example}
\begin{enumerate}[label=(\roman*), leftmargin=2.2em, itemsep=0.15em]
\item 群 $G$ 的表示范畴，张量积取在 $\C$ 上，是编织的，事实上是对称的。
\item 交换代数 $A$ 上的模范畴，张量积取 $\otimes_A$，也是对称的。若 $A$ 不交换，这通常连张量范畴都不是。
\item 如果你没见过非对称的编织范畴，要造一个会更困难。一个著名的例子来自量子群。下面我将给出一个拓扑的例子。
\end{enumerate}
\end{example}

现在给定只有同伦群 $\pi_2,\pi_3$ 的 $X$，令 $\mathcal{B}_X$ 为二重环路空间 $\Omega^2X$ 上、支集在有限多个连通分支上的平坦向量丛范畴。$X$ 的基点的选择无关紧要。于是，我们对 $\pi_2$ 的每个元素选取了 $\pi_3$ 的一个表示。这个范畴是阿贝尔的；若同伦群有限，它甚至是半单的。

记 $m\colon\Omega^2X\times\Omega^2X\to\Omega^2X$ 为乘法，并通过 $U\odot V:=m_*(U\boxtimes V)$（沿 $m$ 的逐纤维零阶同调）在 $\mathcal{B}_X$ 上定义一个张量结构。

\begin{proposition}上述构造在 $\mathcal{B}_X$ 上定义一个编织张量结构。若 $k_3=0$，则该编织是对称的。\end{proposition}

\begin{proof}$\odot$ 的结合性应当清楚；我们只指出编织。关于 $\pi_2$ 交换性的通常论证，利用二重环路空间结构，给出乘法 $m$ 与其转置 $m\circ\tau$ 之间的一个同伦。回忆该构造：把 $\Omega^2X$ 实现为从正方形到 $X$、并把边界送到基点的映射；于是，嵌入一个更大正方形中、表示 $\Omega^2X$ 中乘积的两个相邻正方形 $\square\square$，可以通过顺时针旋转连续地互相越过。这样，我们便用一族单参数映射把 $\odot$ 与 $\odot\circ\tau$ 联结起来。沿这族映射的逐纤维同调丛带有平坦联络；沿区间追随它，便给出编织同构 $\beta(U,V)\colon U\odot V\to V\odot U$。\end{proof}

\begin{example}这里把 $\mathcal{B}_{S^2}$ 写开。回忆 $\pi_2(S^2)\cong\pi_3(S^2)=\Z$，而二次的 Hopf 映射 $\Z\to\Z$ 是平方映射。对象是 $\C[z^{\pm1}]$ 上的 $\Z$-分次模；分次反映丛的支集在 $\pi_2(S^2)$ 中的点，而 $z$ 作为 $\pi_3$ 的生成元作用。乘法于是由在 $\C[z^{\pm1}]$ 上作张量积给出。这自然是对称幺半的，但我们如下修改显然的对称编织：对于支集在次数 $m,n$ 中的模 $U,V$，修改后的编织 $\beta(U,V)\colon U\odot V\to V\odot U$ 是 $z^{mn}$ 的作用。\end{example}

\subsection{带状结构}

不变量 $k_3$ 是 Whitehead 括号的二次加细，所含信息比编织更精细。（因此，即使 $k_3\neq0$，$\mathcal{B}_X$ 仍可能是对称的。）具体地说，$k_3$ 给出 $\mathcal{B}_X$ 上的一个\emph{带状结构}（ribbon structure）：即对每个对象 $x$ 给出一个中心自同构 $\theta_x$，满足
\[
\theta_{x\otimes y}\circ(\theta_x^{-1}\otimes1)\circ(1\otimes\theta_y^{-1})=\beta(y,x)\circ\beta(x,y),\qquad
\theta_{x^{-1}}=\theta_x.
\]

「中心」指的是整个范畴中的态射，意思是 $\theta$ 是 $\mathcal{B}_X$ 上恒等函子的自同构；而所写的第二个恒等式依赖于 $T$ 的群样结构（每个对象的张量可逆性）。\footnote{通常要求在刚性范畴（带内伴随的范畴，见定义 3.5.6）中有 $\theta_x^{\vee}=(\theta_x)^{\vee}$；甚至可以在没有伴随的情形表述这一条件。}在上面的 $S^2$ 例子中，有 $\theta_m=z^{m^2}$。

带状结构 $\theta$ 的第一个方程可以如下理解。编织张量范畴是范畴的 2-范畴中小 2-圆盘运算子（little 2-disk operad）$E_2$ 上的代数对象。带状则是把结构增强为带标架的小 2-圆盘运算子。所谓「带标架 $E_2$ 上的代数」，是某个定向二维 TQFT 的一部分：它只看到圆的并（作为对象），以及带多个输入洞、但只有一个输出的（若干）圆盘（作为态射）。这个部分 TQFT 把 $\mathcal{B}$ 赋给圆周，把由张量单位诱导的函子 $\mathsf{Vect}\to\mathcal{B}$ 赋给那个出射圆盘，把 $m$ 重乘法 $\mathcal{B}^{\otimes m}\to\mathcal{B}$ 赋给 $m$ 洞圆盘。\footnote{说来令人迷惑：无标架的小 2-圆盘运算子所描述的，恰是带标架圆盘的 TQFT 的相应部分。}任何二重环路空间上都有这个带状结构：把命题 3.3.14 展开，可以把它线性化，得到一个带状范畴。

一个密切相关的解释是：$\theta$ 把编织的平方平凡化成该范畴的编织张量自同构。（编织的平方给出乘法函子 $m\colon\mathcal{T}\otimes\mathcal{T}\to\mathcal{T}$ 的一个自同构，它与编织结构相容。）

\begin{remark}$\theta$ 的第二个、齐次的性质，与向三维 TQFT 结构的维数扩张有关。关于图表组合学与拓扑之间在这些定义表面之下暗燃的惊人联系，见书 [6]。\end{remark}

\subsection{有限同伦型}

最后，我们仿照规范理论的例子，说明 §3.1.6 的增强规范理论如何构造为扩张 TQFT。除了语言——它层层叠叠、充满各阶态射——之外，这与 §3.3.3 中有限群的扩张规范理论没有区别。

设 $X$ 是具有有限同伦群的空间，并固定所要的 TQFT 的维数 $D>0$。高于 $D$ 的同伦群不起作用，将被截断。我们分两步寻求从 $\mathsf{Bord}_D$ 到弱线性 $(D-1)$-范畴的 $D$-范畴 $\mathsf{Cat}_{D-1}$ 的一个（弱）函子：
\begin{enumerate}[label=(\roman*), leftmargin=2.2em, itemsep=0.15em]
\item 首先，我们构造函子 $\mathsf{Bord}_D\to\mathsf{Cor}_D$，映入由对应构成的 $D$-范畴，其对象是只有有限多个同伦群、且这些群均为有限群的拓扑空间。这类似于经典场论；物理学家在经典场论上做路径积分以进行量子化。
\item 其次，我们构造量子化函子 $Q\colon\mathsf{Cor}_D\to\mathsf{Cat}_{D-1}$。
\end{enumerate}

第一部分是无歧义的：每个流形 $M$——它表示无论哪一层的态射——都被送到映射空间 $\mathrm{Map}(M;X)$。

然而，由于我们并未选定弱高阶范畴的定义，实际上无法执行第二步：构造的细节依赖于模型。一种构造使用了 $m$-代数（带有相容乘法层的代数）的概念，已在 [11] 中指出。另一种构造可以用 Morrison--Walker [23] 的「Blob 复形」给出。在两种情形中，加上理论的顶层（向量空间之间的线性映射，或数）都是更精细的一步，并且要用到 $X$ 的有限性。不过，让我在此指出两者共有的想法：构造一个空间的 $m$-线性化函子（它是一个线性 $m$-范畴），使同伦纤维积被送到张量积。（零范畴是一个向量空间；线性 $(-1)$-范畴大概就是一个复数，我猜。）

$Y$ 的 0-范畴化是 $Y$ 上局部常值函数的向量空间。1-范畴化是带平坦联络的向量丛范畴。2-范畴化是带平坦联络的线性范畴丛的 2-范畴（线性范畴的局部常值层）。如此等等。

关键观察是：可以定义空间以 $m$-范畴的局部常值层为系数的上同调，并得到一个 $m$-范畴。这让我们能为高阶范畴的丛定义「反向映射」（wrong-way maps），从而把对应转换为函子。\begin{annote}底本此处印作 higer，应为 higher；中文不体现该拼写笔误，译文按正字处理。[来源：上下文推断]\end{annote}对 $X$ 的有限性条件，是为了停留在可对偶对象（dualizable object，下一节）的世界里；这些对象是 TQFT 的相容性所需要的。

§3.3.7 的编织张量范畴是空间 $X$ 的 3-范畴化的一个模型；我们将在最后一节看到，为什么编织张量范畴是一个具有 3-范畴本性的对象：它有一个「模对象的 3-范畴」。
